% ============================================================
%  cover.tex — 概率论笔记封面（九层设计）
%  参考 Charlie's LaTeX Template 封面架构
%  配色：GreenMain/GreenDark/GreenAccent 绿系 + Gold 暖金
%  用法：\coverimage{figure/封面.png} → \makecover
% ============================================================

% ── 封面图片命令 ──
\makeatletter
\providecommand{\@coverimage}{}
\newcommand{\coverimage}[1]{\renewcommand{\@coverimage}{#1}}

\newcommand{\makecover}{%
\begin{titlepage}
\thispagestyle{empty}

% ============================================================
%  第一层：左侧深绿色竖条 + 渐变纹理 + 装饰圆点 + 细线
% ============================================================
\begin{tikzpicture}[remember picture,overlay]
  % 主体竖条
  \fill[GreenDark!95]
    (current page.south west) rectangle
    ([xshift=2.8cm]current page.north west);
  % 顶部向下的渐变淡出
  \shade[top color=GreenDark!90, bottom color=GreenDark!95]
    ([xshift=0cm,yshift=-4cm]current page.north west) rectangle
    ([xshift=2.8cm]current page.north west);
  % 底部向上的渐变淡出
  \shade[bottom color=GreenDark!90, top color=GreenDark!95]
    ([xshift=0cm,yshift=4cm]current page.south west) rectangle
    ([xshift=2.8cm]current page.south west);
  % 竖条右边界线
  \draw[GreenAccent!50,line width=1pt]
    ([xshift=2.8cm]current page.south west) --
    ([xshift=2.8cm]current page.north west);
  % 左侧装饰细线（距左边缘 0.4cm）
  \draw[GreenLight!25,line width=0.4pt]
    ([xshift=0.4cm]current page.south west) --
    ([xshift=0.4cm]current page.north west);

  % 绿色区域内的装饰圆
  \foreach \y/\r/\o in {
    0.08/1.1/0.07, 0.22/0.7/0.11, 0.40/1.4/0.06,
    0.55/0.55/0.10, 0.68/0.9/0.08, 0.82/0.6/0.11, 0.93/0.75/0.09
  } {
    \fill[white, opacity=\o]
      ([xshift=1.4cm, yshift=-\y\paperheight]current page.north west)
      circle (\r cm);
  }

  % 两条细竖线装饰
  \draw[white, line width=0.3pt, opacity=0.18]
    ([xshift=0.5cm]current page.north west) --
    ([xshift=0.5cm]current page.south west);
  \draw[white, line width=0.3pt, opacity=0.18]
    ([xshift=2.1cm]current page.north west) --
    ([xshift=2.1cm]current page.south west);
\end{tikzpicture}

% ============================================================
%  第二层：顶部/底部装饰金线带 + 等距标记圆点
% ============================================================
\begin{tikzpicture}[remember picture,overlay]
  % 顶部横贯金线
  \draw[GoldMain!50,line width=0.6pt]
    ([xshift=3.0cm,yshift=-1.6cm]current page.north west) --
    ([xshift=-2.0cm,yshift=-1.6cm]current page.north east);
  % 金线上方小标记点
  \foreach \x in {3.5,7,10.5,14} {
    \fill[GoldMain!60]
      ([xshift=\x cm,yshift=-1.45cm]current page.north west)
      circle (1.5pt);
  }
  % 底部横贯金线
  \draw[GoldMain!40,line width=0.5pt]
    ([xshift=3.0cm,yshift=2.2cm]current page.south west) --
    ([xshift=-1.5cm,yshift=2.2cm]current page.south east);
  % 底部金线标记
  \foreach \x in {4,8,12} {
    \fill[GoldMain!45]
      ([xshift=\x cm,yshift=2.35cm]current page.south west)
      circle (1.2pt);
  }
\end{tikzpicture}

% ============================================================
%  第三层：右上角装饰几何图形组（同心圆环 + 放射线）
% ============================================================
\begin{tikzpicture}[remember picture,overlay]
  % 大圆环（空心）
  \draw[GoldMain!30,line width=0.7pt]
    ([xshift=-2.5cm,yshift=-2.5cm]current page.north east)
    circle (22pt);
  % 小圆环（同心）
  \draw[GreenAccent!30,line width=0.5pt]
    ([xshift=-2.5cm,yshift=-2.5cm]current page.north east)
    circle (14pt);
  % 中心实心点
  \fill[GoldMain!50]
    ([xshift=-2.5cm,yshift=-2.5cm]current page.north east)
    circle (2.5pt);
  % 正交放射线
  \foreach \ang in {0,90,180,270} {
    \draw[GoldMain!25,line width=0.4pt]
      ([xshift=-2.5cm,yshift=-2.5cm]current page.north east)
      -- ++(\ang:28pt);
  }
  % 对角放射线
  \foreach \ang in {45,135,225,315} {
    \draw[GreenAccent!20,line width=0.3pt]
      ([xshift=-2.5cm,yshift=-2.5cm]current page.north east)
      -- ++(\ang:26pt);
  }
\end{tikzpicture}

% ============================================================
%  第四层：左下角装饰 —— 概率分布曲线簇（正态型 + 指数型）
% ============================================================
\begin{tikzpicture}[remember picture,overlay]
  % 正态分布型曲线（钟形）
  \draw[GoldMain!15,line width=0.5pt,smooth]
    ([xshift=3.5cm,yshift=3.0cm]current page.south west)
    to[out=30,in=210] ++(3.0,2.0);
  \draw[GreenAccent!15,line width=0.5pt,smooth]
    ([xshift=3.3cm,yshift=3.0cm]current page.south west)
    to[out=35,in=205] ++(3.2,2.2);
  \draw[GoldMain!12,line width=0.5pt,smooth]
    ([xshift=3.7cm,yshift=2.8cm]current page.south west)
    to[out=25,in=215] ++(2.8,1.8);
  % 指数衰减型曲线
  \draw[GreenMain!12,line width=0.5pt,smooth]
    ([xshift=3.5cm,yshift=4.0cm]current page.south west)
    to[out=-15,in=195] ++(2.5,-1.5);
  \draw[GoldMain!10,line width=0.4pt,smooth]
    ([xshift=3.7cm,yshift=4.2cm]current page.south west)
    to[out=-20,in=190] ++(2.3,-1.3);
  \draw[GreenAccent!10,line width=0.4pt,smooth]
    ([xshift=3.9cm,yshift=3.8cm]current page.south west)
    to[out=-10,in=200] ++(2.6,-1.0);
  % 流方向小箭头
  \draw[GoldMain!25,-{Stealth[length=3pt,width=3pt]}]
    ([xshift=5.0cm,yshift=4.2cm]current page.south west)
    -- ++(0.8,0.5);
  \draw[GreenAccent!20,-{Stealth[length=3pt,width=3pt]}]
    ([xshift=5.2cm,yshift=3.5cm]current page.south west)
    -- ++(0.7,-0.5);
\end{tikzpicture}

% ============================================================
%  第五层：公式背景区 —— 淡金色矩形 + 概率论经典公式
% ============================================================
\begin{tikzpicture}[remember picture,overlay]
  % 主背景矩形（含极淡渐变）
  \shade[top color=GoldPale!40, bottom color=GoldPale!25]
    ([xshift=3.3cm,yshift=-3cm]current page.north west) rectangle
    ([xshift=-1.5cm,yshift=2.5cm]current page.south east);
  % 矩形的细金边框
  \draw[GoldMain!30,line width=0.4pt]
    ([xshift=3.3cm,yshift=-3cm]current page.north west) rectangle
    ([xshift=-1.5cm,yshift=2.5cm]current page.south east);
  % 公式间递进连接箭头
  \draw[GoldMain!35,line width=0.4pt,-{Stealth[length=3.5pt,width=3.5pt]}]
    ([xshift=5.0cm,yshift=-4.3cm]current page.north west) --
    ([xshift=5.0cm,yshift=-4.8cm]current page.north west);
  \draw[GoldMain!35,line width=0.4pt,-{Stealth[length=3.5pt,width=3.5pt]}]
    ([xshift=5.2cm,yshift=-5.3cm]current page.north west) --
    ([xshift=5.2cm,yshift=-5.8cm]current page.north west);

  % 六条概率论经典公式
  \node[anchor=west,GreenMain!70] at ([xshift=3.7cm,yshift=-3.5cm]current page.north west)
    {$\displaystyle \PP\!\left(\bigcup_{i=1}^{\infty}A_i\right)=\sum_{i=1}^{\infty}\PP(A_i)$};

  \node[anchor=west,GreenAccent!65] at ([xshift=3.8cm,yshift=-4.5cm]current page.north west)
    {$\displaystyle \PP(A\mid B)=\frac{\PP(AB)}{\PP(B)}$};

  \node[anchor=west,GoldMain!75] at ([xshift=3.9cm,yshift=-5.5cm]current page.north west)
    {$\displaystyle \EE[X]=\int_{\Omega}X\,\dd\PP$};

  \node[anchor=west,GreenMain!55] at ([xshift=3.7cm,yshift=-6.5cm]current page.north west)
    {$\displaystyle X_n\xrightarrow{\;\PP\;}X\;\Longleftrightarrow\; \forall\varepsilon>0,\; \PP(|X_n-X|>\varepsilon)\to0$};

  \node[anchor=west,GreenAccent!50] at ([xshift=3.8cm,yshift=-7.5cm]current page.north west)
    {$\displaystyle \frac{1}{n}\sum_{i=1}^{n}X_i\xrightarrow{\text{a.s.}}\EE[X_1]$};

  \node[anchor=west,GoldMain!55] at ([xshift=4.0cm,yshift=-8.5cm]current page.north west)
    {$\displaystyle f_{X}(x)=\frac{1}{\sqrt{2\pi\sigma^{2}}}\,
       \exp\!\left(-\frac{(x-\mu)^{2}}{2\sigma^{2}}\right)$};

  % 公式编号装饰（右侧小圆点）
  \foreach \y/\c in {-3.5/GreenMain!35, -4.5/GreenAccent!30, -5.5/GoldMain!35, -6.5/GreenMain!30, -7.5/GreenAccent!25, -8.5/GoldMain!30} {
    \fill[\c]
      ([xshift=-2.0cm,yshift=\y cm]current page.north west)
      circle (1.8pt);
  }
\end{tikzpicture}

% ============================================================
%  第六层：标题区 —— 装饰分隔线 + 标题文字
% ============================================================
\vspace*{2cm}

\hspace*{3.3cm}
\begin{minipage}{0.7\textwidth}
  % ── 系列标识 ──
  {\large\bfseries\color{GreenAccent} Huginn Notes in Mathematics\par}

  \vspace{0.3cm}
  % ── 装饰短金线（带渐小圆点尾饰）──
  \begin{tikzpicture}
    \draw[GoldMain!60,line width=1.2pt] (0,0) -- (2.0,0);
    \fill[GoldMain!70] (2.3,0) circle (2pt);
    \fill[GoldMain!50] (2.8,0) circle (1.2pt);
    \fill[GoldMain!30] (3.2,0) circle (0.8pt);
  \end{tikzpicture}

  \vspace{0.8cm}

  % ── 主标题 ──
  {\Huge\bfseries\color{GreenDark} \@title\par}

  \vspace{0.6cm}
  % ── 英文副标题 ──
  {\Large\color{GreenMain!70} \@author\par}

  \vspace{0.4cm}
  % ── 第二道装饰金线（两端圆点）──
  \begin{tikzpicture}
    \draw[GoldMain!40,line width=0.7pt] (0,0) -- (3.5,0);
    \fill[GoldMain!45] (0,0) circle (1.5pt);
    \fill[GoldMain!45] (3.5,0) circle (1.5pt);
  \end{tikzpicture}

  \vspace{0.8cm}
  % ── Fraktur 铭文 ──
  {\large\color{GoldMain!85} $\mathfrak{For}\quad\mathfrak{the}\quad\mathfrak{Worthy}$\par}

  \vspace{1.0cm}


\end{minipage}

% ============================================================
%  第七层：右下角「出版社」小 logo（仿 Springer 风格）
% ============================================================
\begin{tikzpicture}[remember picture, overlay]
  \begin{scope}[shift={([xshift=-5.5cm, yshift=1.5cm]current page.south east)}]
    % 小房子
    \draw[GreenMain, line width=0.8pt] (0,0) rectangle (0.9,0.72);
    \draw[GreenMain, line width=0.8pt] (0,0.72) -- (0.45,1.08) -- (0.9,0.72);
    \draw[GoldMain, line width=0.5pt] (0.45,0.72) -- (0.45,1.08);
    \draw[GreenMain, line width=0.4pt] (0.23,0) -- (0.23,0.72);
    \draw[GreenMain, line width=0.4pt] (0.67,0) -- (0.67,0.72);
    % 屋顶金点
    \fill[GoldMain] (0.45,1.08) circle (1.5pt);
    \node[anchor=west, GreenMain, font=\footnotesize] at (1.1,0.32)
      {Produced by Huginn};
  \end{scope}
\end{tikzpicture}

% ============================================================
%  第八层：底部主封面图片 —— 与金色框左右对齐，位于小房子上方
% ============================================================
\ifdefempty{\@coverimage}{}{
	\begin{tikzpicture}[remember picture,overlay]
		\node[
		anchor=south west,
		inner sep=0pt,
		outer sep=0pt
		] at ([xshift=3.3cm,yshift=3.05cm]current page.south west)
		{
			\includegraphics[
			width=\dimexpr\paperwidth-4.8cm\relax
			]{\@coverimage}
		};
	\end{tikzpicture}
}

% ============================================================
%  第九层：散落装饰点（页面随机分布的数学点缀）
% ============================================================
\begin{tikzpicture}[remember picture,overlay]
  % 右中区域散点
  \fill[GoldMain!20] ([xshift=-3.0cm,yshift=-9.5cm]current page.north west) circle (1.0pt);
  \fill[GreenAccent!20] ([xshift=-3.8cm,yshift=-9.8cm]current page.north west) circle (0.7pt);
  \fill[GoldMain!18] ([xshift=-2.5cm,yshift=-10.2cm]current page.north west) circle (0.8pt);
  % 右下方散点
  \fill[GreenMain!15] ([xshift=-4.0cm,yshift=3.5cm]current page.south west) circle (1.2pt);
  \fill[GoldMain!15] ([xshift=-5.5cm,yshift=4.0cm]current page.south west) circle (0.9pt);
  \fill[GreenAccent!12] ([xshift=-4.8cm,yshift=3.0cm]current page.south west) circle (0.6pt);
  % 标题区附近轻微装饰
  \fill[GoldMain!22] ([xshift=14.5cm,yshift=-6.0cm]current page.north west) circle (1.0pt);
  \fill[GreenAccent!18] ([xshift=14.0cm,yshift=-5.5cm]current page.north west) circle (0.6pt);
  \fill[GoldMain!18] ([xshift=15.0cm,yshift=-5.8cm]current page.north west) circle (0.8pt);
\end{tikzpicture}

\end{titlepage}
}
\makeatother
